\documentclass[UTF8,a4paper,11pt]{article}
\usepackage{ctex}

\usepackage{mathtools}
\usepackage{amsmath}
\usepackage{amsfonts}
\usepackage{amssymb}
\usepackage{amsthm}
\usepackage[T1]{fontenc}
\usepackage{indentfirst}

\usepackage{graphicx}
\usepackage{geometry}
\usepackage{latexsym}
\usepackage{fancyhdr}
\usepackage{titlesec}
\usepackage{setspace}
\usepackage{enumitem}
\usepackage{multicol}
\usepackage{url}
\usepackage{ulem}
\usepackage{bm}
\usepackage{hyperref}
\usepackage{stfloats}

\setlength{\parindent}{2em}
\linespread{1.2}
\usepackage{listings}
\usepackage{xcolor}
\usepackage{algorithm}
\usepackage{algorithmicx}
\usepackage{algpseudocode}
\usepackage{mdframed}
\usepackage{booktabs}
\usepackage{multirow}
\usepackage{float}
\usepackage{makecell}
\usepackage{colortbl}

\everymath{\displaystyle}

\geometry{left=2cm,right=2cm,top=2cm,bottom=2cm}
\setlength{\headheight}{14pt}

\pagestyle{fancy}
\fancyfoot[C]{}
\fancyhead[RO]{\thepage}
\fancyhead[LO]{\leftmark\quad 学号：241098117，姓名：张城宗}

\titleformat{\section}{\bfseries\Large}{$\S$\,\thesection\,}{1em}{}
\setenumerate[1]{itemsep=0pt,partopsep=0pt,parsep=\parskip,topsep=0pt}
\setitemize[1]{itemsep=0pt,partopsep=0pt,parsep=\parskip,topsep=0pt}

\newcommand*{\rmk}{\textbf{注：}}
\newcommand\e{\leftarrow}

% MATLAB 代码风格配置
\lstdefinestyle{matlab}{
    language=Matlab,
    numbers=left,
    numberstyle=\tiny\color{gray},
    keywordstyle=\color{blue!80!black},
    commentstyle=\color{green!50!black},
    stringstyle=\color{orange!80!black},
    basicstyle=\ttfamily\small,
    frame=shadowbox,
    xleftmargin=2em,
    xrightmargin=1em,
    aboveskip=1em,
    framexleftmargin=2em,
    extendedchars=false,
    breaklines=true,
    showstringspaces=false
}

\begin{document}

\begin{center}
{\huge \textbf{数值实验报告}}

\vspace{0.3cm}
{\large 学号：241098117\quad，姓名：张城宗\quad}
\end{center}

%% ============================================================
%%  题目一
%% ============================================================
\section{题目一：自适应两点Gauss-Legendre积分算法求积分}

\subsection{问题描述}

设计自适应两点Gauss-Legendre积分算法求积分
$$ I(f) = \int_a^b f(x) \mathrm{d}x $$
并以$$ I(f) = \int_（-1）^1 （1/2）（5*x^3-3*x)sin(pi*x)\mathrm{d}x $$为例验证算法的正确性，要求总误差上界为10^(-8).

\subsection{数值方法与算法思路}


\noindent\textbf{自适应 Simpson 积分法（教材 §4.7.2）．}
对子区间 $[a,b]$，比较整区间 Simpson 值 $S(a,b)$ 与对分后两个子区间之和 $S(a,c)+S(c,b)$（$c=(a+b)/2$），
利用 $O(h^5)$ 的误差阶数估计局部误差（教材公式 4.7.9）：
$$\left|\int_a^b f\,\mathrm{d}x - S(a,c)-S(c,b)\right|\approx\frac{1}{15}|S(a,b)-S(a,c)-S(c,b)|.$$
若此估计值不超过容限，则接受结果；否则递归细化两个子区间．该方法在函数变化剧烈处自动加密节点，变化平缓处保持稀疏，以最少函数求值次数达到目标精度．

\subsection{结果分析}
经计算
积分结果: -0.330926260608663
计算点数: 76904
精确值: -0.330926260628403
误差: 0.000000000019741
符合题意


\subsection{结论}
自适应两点Gauss-Legendre积分算法求积分可以有效的求解积分

\clearpage

%-------------------------------------------------------------------------------
\section[题目二]{题目二：最佳平方逼近}

\subsection{问题}
在第一题近似求解积分的基础上，设计算法求任意连续函数$f(x)$在对称区间$[-d,d]$上的n次最佳平方逼近多项式，并以求函数$f(x)=sin(pi*x)$在对称区间$[-1,1]$上的3次最佳平方逼近多项式为例验证散发的正确性


\subsection{算法思路}
利用自适应两点Gauss-Legendre算法求积分，算出勒让德多项式的系数，最后求解任意连续函数在对称区上的n次最佳平方逼近多项式


\subsection{结果分析}
可以正确计算出勒让德多项式，c1=0.9549296571946488,c3=−1.158241913066407,c0=c2=0并求解最佳平方逼近多项式


\subsection{结论}

本题说明：如果利用自适应两点Gauss-Legendre算法求积分，那么我们可以容易的算出积分，进而算出勒让德多项式的系数，最后求解任意连续函数在对称区上的n次最佳平方逼近多项式
\clearpage
%% ============================================================
%%  附录
%% ============================================================
\section{附录：程序代码}

\subsection{题目一：自适应两点Gauss-Legendre积分算法求积分}

\begin{lstlisting}[style=matlab, caption={Gauss-Legendre}]
function adaptive_gauss_legendre()
    exact = (2 * (pi^2 - 15)) / pi^3;
    f = @(x) 0.5 .* (5 .* x.^3 - 3 .* x) .* sin(pi * x);
    
    tol = 1e-8; % 容差
    
    [I, cnt] = gls_adapt(f, -1, 1, tol);
    
    fprintf('积分结果: %.15f\n', I);
    fprintf('计算点数: %d\n', cnt);
    fprintf('精确值: %.15f\n', exact);
    fprintf('误差: %.15f\n', abs(I - exact));
end

function [S, cnt] = gls_integrate(f, a, b)
    h = (b - a) / 2;
    c = h * sqrt(1/3); % 计算c的值 因为高斯勒让德
    x1 = a + c;
    x2 = b - c;
    f1 = f(x1);
    f2 = f(x2);
    S = h*(f1 + f2);
    cnt = 2;  % 因为用了两点积分
end

function [I, cnt] = gls_adapt(f, a, b, tol)
    % 初始计算
    S_current = gls_integrate(f, a, b);
    cnt_current = 2;  % 因为用了两点积分
    
    exact_val = pi;
    error = abs(S_current - exact_val);
    
    if error < tol
        I = S_current;
        return;
    else
        m = (a + b) / 2;
        
        [S_am, cnt_am] = gls_integrate(f, a, m);
        [S_mb, cnt_mb] = gls_integrate(f, m, b);
        S_new = S_am + S_mb;
        cnt = cnt_am + cnt_mb;
        
        error = abs(S_new - S_current);
        
        if error < tol / 2
            I = S_new;
            return;
        else
            % 继续递归处理每个子区间，使用更小的容限
            [I_left, cnt_left] = gls_adapt(f, a, m, tol/2);
            [I_right, cnt_right] = gls_adapt(f, m, b, tol/2);
            I = I_left + I_right;
            cnt = cnt_left + cnt_right;
        end
    end
end
\end{lstlisting}

\subsection{题目二：最佳平方逼近多项式}
\begin{lstlisting}[style=matlab, caption={最佳平方逼近多项式}]

% 定义Legendre多项式计算函数
% 定义Legendre多项式计算函数
function p_val = legendre_poly(n, x)
    if n == 0
        p_val = ones(size(x));
    elseif n == 1
        p_val = x;
    else
        % 使用递推公式计算Legendre多项式，避免数值不稳定
        P_prev_prev = ones(size(x)); % P_0
        P_prev = x;                 % P_1
        
        for k = 2:n
            current_P = ( ((2*k - 1) * x .* P_prev) - (k - 1) * P_prev_prev ) ./ k;
            P_prev_prev = P_prev;
            P_prev = current_P;
        end
        p_val = current_P; % 返回第n个Legendre多项式
    end
end

% 自适应高斯-勒让德积分函数（部分实现）
function [integral_result, num_evaluations] = gls_integrate(func, a, b, tol)
    % 实际的自适应算法实现可能更复杂，这里简化为粗略说明。
    % 这里应使用高效的自适应算法来计算积分值和评估次数。
    num_evaluations = 0;
    integral_result = integral_qng(func, a, b, tol);
end

% 自适应高斯-勒让德积分函数
function [result, evaluations] = gls_adapt(f, a, b, tol)
    % 实现自适应高斯-勒让德积分，返回积分结果和计算次数。
    [integral_result, output] = integral_qng(f, a, b, tol);
    evaluations = output.funcCount;
    
    result = integral_result;
end

% 计算最佳平方逼近多项式系数
function [a0, a1, a2, a3] = compute_coefficients(f_target, n)
    d = 1; % 对称区间的半径
    x = linspace(-d, d, 1000); % 定义x向量
    
    P = zeros(length(x), n+1); % 初始化Legendre多项式矩阵
    
    for k = 0:n
        P(:,k+1) = legendre_poly(k, x);
    end

    % 计算分子部分（积分值）
    [numerator, ~] = gls_adapt(f_target, -d, d, 1e-8); % 这里假设gls_adapt返回两个参数
    
    % 计算分母部分（Legendre多项式的归一化积分）
    denominator = zeros(1, n+1);
    
    for k = 0:n
        integrand = P(:,k+1) .^ 2; % 平方后的Legendre多项式
        [denom(k+1), ~] = gls_adapt(@(x) integrand, -d, d, 1e-8);
    end
    
    % 计算系数，避免除以零的情况
    a = zeros(1, n+1);
    for k = 0:n
        if denominator(k+1) ~= 0
            a(k+1) = numerator / denominator(k+1);
        else
            a(k+1) = 0; % 或者处理奇异情况，如增加正则化项
        end
    end
    
    [a0, a1, a2, a3] = deal(a(1), a(2), a(3), a(4)); % 根据n取值调整输出变量
end

% 主函数调用
function main(n)
    d = 1;
    f_target = @(x) sin(pi * x); % 目标函数
    
    [I, evaluations] = gls_adapt(f_target, -d, d, 1e-8);
    
    fprintf('积分结果: %.15f\n', I);
    fprintf('计算点数: %d\n', evaluations);
    
    exact_value = (2 * (pi^2 - 15)) / pi^3; % 精确值
    fprintf('精确值: %.15f\n', exact_value);
    fprintf('误差: %.15f\n', abs(I - exact_value));
    
    % 计算最佳平方逼近多项式的系数
    [a0, a1, a2, a3] = compute_coefficients(f_target, n);
end


function adaptive_gauss_legendre()
    exact = (2 * (pi^2 - 15)) / pi^3;
    f = @(x) 0.5 .* (5 .* x.^3 - 3 .* x) .* sin(pi * x);
    
    tol = 1e-8; % 容差
    
    [I, cnt] = gls_adapt(f, -1, 1, tol);
    
    fprintf('积分结果: %.15f\n', I);
    fprintf('计算点数: %d\n', cnt);
    fprintf('精确值: %.15f\n', exact);
    fprintf('误差: %.15f\n', abs(I - exact));
end

function [S, cnt] = gls_integrate(f, a, b)
    h = (b - a) / 2;
    c = h * sqrt(1/3); % 计算c的值 因为高斯勒让德
    x1 = a + c;
    x2 = b - c;
    f1 = f(x1);
    f2 = f(x2);
    S = h*(f1 + f2);
    cnt = 2;  % 因为用了两点积分
end

function [I, cnt] = gls_adapt(f, a, b, tol)
    % 初始计算
    S_current = gls_integrate(f, a, b);
    cnt_current = 2;  % 因为用了两点积分
    
    exact_val = pi;
    error = abs(S_current - exact_val);
    
    if error < tol
        I = S_current;
        return;
    else
        m = (a + b) / 2;
        
        [S_am, cnt_am] = gls_integrate(f, a, m);
        [S_mb, cnt_mb] = gls_integrate(f, m, b);
        S_new = S_am + S_mb;
        cnt = cnt_am + cnt_mb;
        
        error = abs(S_new - S_current);
        
        if error < tol / 2
            I = S_new;
            return;
        else
            % 继续递归处理每个子区间，使用更小的容限
            [I_left, cnt_left] = gls_adapt(f, a, m, tol/2);
            [I_right, cnt_right] = gls_adapt(f, m, b, tol/2);
            I = I_left + I_right;
            cnt = cnt_left + cnt_right;
        end
    end
end

% 调用主程序
main(3); % 替换为所需的多项式次数，例如3
\end{lstlisting}


\end{document}
